\chapter{نتیجه‌گیری و پیشنهادها}
\label{ch:conclusion}

\section{خلاصه و دستاوردها}
در این پژوهش، معماری \lr{DuoDiT} به عنوان یک رویکرد دوجریانی برای تنظیم دقیق پارامتر-کارای مدل‌های انتشار مبتنی بر مبدل معرفی شد. هسته اصلی روش، حفظ بدنه پیش‌آموزش‌دیده \lr{DiT-XL/2} به صورت منجمد و افزودن یک شاخه دوم سبک با وضوح توکنی بالاتر است. این شاخه از طریق توکن‌های کلاس محلی با جریان اصلی هم‌تراز می‌شود و بدون آموزش مجدد کامل بدنه، امکان تزریق ویژگی‌های محلی و ریزدانه را فراهم می‌کند \cite{2212.09748,8409325}.

مهم‌ترین دستاوردهای این پژوهش عبارتند از:
\begin{enumerate}
    \item \textbf{معماری دوجریانی \lr{DuoDiT}}: طراحی یک شاخه دوم با اندازه پچ کوچک‌تر که در کنار بدنه اصلی منجمد عمل می‌کند و ویژگی‌های محلی را پیش از ادغام با جریان اصلی استخراج می‌نماید.
    \item \textbf{مکانیزم توکن کلاس محلی}: یک راهبرد محتوا-آگاه برای تجمیع ویژگی‌های شاخه دوم که در مطالعه حذفی نسبت به میانگین‌گیری، بیشینه‌گیری و پروجکشن خطی عملکرد بهتری ثبت کرده است.
    \item \textbf{ارزیابی بزرگ‌مقیاس}: ارزیابی روی \lr{ImageNet-1K} نشان داد که \lr{DuoDiT-200} با \lr{17.95M} پارامتر آموزش‌پذیر از مجموع \lr{690.39M} پارامتر، به \lr{FID=2.219}، \lr{KID=0.0014} و \lr{LPIPS-Mean=0.7162} می‌رسد \cite{heusel2017fid,binkowski2018kid,zhang2018lpips}.
    \item \textbf{شواهد ادراکی انسانی}: در مطالعه ترجیح انسانی میان \lr{DuoDiT} و \lr{LightningDiT}، مدل پیشنهادی در \lr{53.26\%} مقایسه‌های غیرمساوی ترجیح داده شد، که با بهبود \lr{LPIPS-Mean} همسو است \cite{Yao2025LightningDiT}.
\end{enumerate}

\section{پیشنهادهایی برای پژوهش‌های آینده}
بر اساس یافته‌های این پژوهش، مسیرهای زیر برای تحقیقات آتی پیشنهاد می‌شود:
\begin{enumerate}
    \item \textbf{بررسی شرط‌گذاری اختیاری \lr{FiLM}}: فعال‌سازی و ارزیابی کنترل‌شده \lr{FiLM} در تنظیمات شرطی ریزدانه‌تر می‌تواند نشان دهد آیا همگام‌سازی مستقیم شاخه دوم با زمان و کلاس بهبود معناداری ایجاد می‌کند یا خیر \cite{perez2018film}.
    \item \textbf{شاخه‌های کمکی چندمقیاسی}: توسعه بیش از یک شاخه وضوح‌بالا می‌تواند به مدل اجازه دهد جزئیات محلی را در چند مقیاس مکانی همزمان پردازش کند.
    \item \textbf{کاهش هزینه استنتاج}: ترکیب \lr{DuoDiT} با روش‌هایی مانند کوانتیزاسیون پس از آموزش یا مسیرهای محاسباتی پویا می‌تواند هزینه اجرای مدل را کاهش دهد \cite{2406.17343}.
    \item \textbf{تعمیم به تولید متن-به-تصویر و ویدئو}: استفاده از شاخه دوم در مدل‌های شرطی متنی یا ویدئویی می‌تواند نشان دهد آیا این ایده فراتر از تولید کلاسی \lr{ImageNet} نیز مفید است یا خیر \cite{9851474,8313484}.
    \item \textbf{جایگزینی بدنه اصلی/خبره}: جایگزینی اجزای اصلی مدل، مانند بدنه اصلی یا خبره، با مدل‌های دیگر می‌تواند منجر به بهبود عملکرد و کاهش هزینه محاسباتی شود.
\end{enumerate}

در مجموع، \lr{DuoDiT} نشان می‌دهد که اصلاحات معماری هدفمند می‌توانند مسیر مؤثری برای تنظیم دقیق پارامتر-کارای مدل‌های انتشار مبتنی بر مبدل باشند. این روش جایگزین آموزش بزرگ‌مقیاس مدل‌های قدرتمندتر نیست، بلکه راهکاری برای بهره‌گیری از دانش یک بدنه منجمد و افزودن ظرفیت محلی با هزینه محدود فراهم می‌کند.
